\section{Study Plant and Dataset}
\label{sec:dataset}

The study uses a \Cadence{} operational dataset for a \PlantCapacity{} PV plant at \StudySite{} (\StudyCoords). The records combine plant-monitoring values with time-aligned meteorological, cloud, and resource channels on the same IST timestamp grid. The active dataset spans \DataSpan{} and contains 525,888 rows after the data-cleaning phase. Machine-fault periods, hardware-level errors, plant shutdown intervals, and invalid sensor states were omitted before model training and evaluation. Weather-driven variability is retained, including cloudy intervals, monsoon-like conditions, high humidity, rain, and low-irradiance periods.

\begin{table}[!t]
\caption{Study plant dataset summary.}
\label{tab:dataset_summary}
\centering
\small
\begin{tabularx}{\columnwidth}{@{}l>{\raggedright\arraybackslash}X@{}}
\toprule
Property & Value \\
\midrule
Plant rating & 43.624 kW \\
Location & VIT Chennai, Tamil Nadu \\
Coordinates & 12.8419815 N, 80.1549340 E \\
Time span & 2021-01-01 to 2025-12-31 \\
Sampling interval & 5 min \\
Data source & Operational plant logs with aligned meteorological, cloud, and resource channels \\
Rows & 525,888 \\
Recorded variables & 12 \\
Main target & AC plant power (kW) \\
Resource channel & measured irradiance \\
Meteorological variables & temperature, humidity, wind speed, precipitation \\
Cloud variables & total, low, mid, and high cloud cover \\
Maximum observed power & 43.624 kW \\
\bottomrule
\end{tabularx}
\end{table}


The raw variables used by the forecasting workflow are time, irradiance, power, theoretical power, temperature, humidity, wind speed, precipitation, total cloud cover, low cloud cover, mid cloud cover, and high cloud cover. The irradiance column is treated as a measured resource channel, not as a meteorological input available to the proposed deployable path at the target timestamp. The theoretical-power channel provides a physical reference curve for the plant envelope. The cloud-layer variables are kept separate because different cloud heights can influence irradiance differently. Low clouds often correspond to sharp attenuation and rapid ramps, whereas high cloud cover can produce smoother irradiance reductions.

\begin{table*}[!t]
\caption{Raw dataset variables and their role in the forecasting pipeline.}
\label{tab:raw_variable_roles}
\centering
\small
\begin{tabularx}{\textwidth}{@{}l l >{\raggedright\arraybackslash}X@{}}
\toprule
Variable & Unit/type & Manuscript role \\
\midrule
\texttt{time} & timestamp & Defines the 5-min grid, chronological split, calendar harmonics, lags, and rolling windows. \\
\texttt{irradiance} & W/m$^2$ & Measured resource channel used to construct the Stage-1 $\kt$ target and the sensor-assisted diagnostic model; not used as a measured target-time input in the proposed forecast path. \\
\texttt{power} & kW & Main target variable for plant-level AC power forecasting and all persistence baselines. \\
\texttt{theoretical\_power} & kW & Logged theoretical-power reference channel used for interpretation and the sensor-assisted diagnostic path; the proposed path instead uses $P_{\mathrm{cs}}(t)$ and $\widehat{P}_{\mathrm{theo},c}(t)$ derived from timestamp, location, and Stage-1 estimates. \\
\texttt{temp} & $^\circ$C & Meteorological input; contributes to seasonal behavior, dew-point approximation, and temperature-related atmospheric variation. \\
\texttt{humidity} & \% & Meteorological input; used directly, as monthly anomaly, and in dew-point/spread features. \\
\texttt{wind\_speed} & m/s & Meteorological input; proxy for local atmospheric movement and cooling/ramp context. \\
\texttt{precipitation} & mm & Meteorological input; rolling sums capture rain and storm context over 1-h, 3-h, and 24-h windows. \\
\texttt{cloud\_cover} & \% & Total cloud-cover predictor and source for lag, rolling, and change features. \\
\texttt{cloud\_cover\_low} & \% & Cloud-layer predictor associated with rapid attenuation and sharp ramps. \\
\texttt{cloud\_cover\_mid} & \% & Cloud-layer predictor that helps separate mid-level attenuation from total cloudiness. \\
\texttt{cloud\_cover\_high} & \% & Cloud-layer predictor associated with smoother high-cloud attenuation and diffuse-light effects. \\
\bottomrule
\end{tabularx}
\end{table*}


Measured and estimated irradiance play different roles. Measured irradiance is available in the historical dataset and is used to construct the Stage-1 target and the sensor-assisted diagnostic model. The proposed deployable path, however, does not assume that measured irradiance at the target timestamp is available. It estimates a transmissivity state from meteorological and cloud variables, then converts that state into an irradiance proxy and a theoretical-power envelope. The sensor-assisted model is therefore reported as a diagnostic upper bound rather than as the headline proposed-path result.

The maximum recorded plant power is 43.624 kW, matching the nominal plant capacity used in the project context. The mean plant power is 9.329 kW across all timestamps, including night and low-sun intervals. Irradiance ranges from 0 to 1125 W/m$^2$, while theoretical power reaches 56.439 kW under idealized clear-sky-like conditions. This separation between theoretical and measured power lets the learning problem represent atmospheric transmissivity and plant conversion behavior rather than simply fitting a repeating daily curve.

\begin{table}[!t]
\caption{Dataset quality and split audit used for manuscript evaluation.}
\label{tab:data_quality_summary}
\centering
\small
\begin{tabularx}{\columnwidth}{@{}l>{\raggedright\arraybackslash}X@{}}
\toprule
Check & Result \\
\midrule
Time span & 2021-01-01 00:00:00 to 2025-12-31 23:55:00 \\
Rows after timestamp de-duplication & 525,888 \\
Duplicate timestamps removed & 0 \\
Dominant cadence & 5.0 min \\
Detected cadence gaps & 0 \\
Missing values & none \\
Machine/hardware fault handling & Fault, shutdown, and invalid sensor periods omitted during data cleaning \\
Power range & 0.000 to 43.624 kW \\
Irradiance range & 0.0 to 1125.0 W/m$^2$ \\
Raw train/validation/test rows & 315,360 / 104,256 / 106,272 \\
Feature-available train/validation/test rows & 315,348 / 104,256 / 106,272 \\
Raw rows in test interval & 106,272 \\
Daylight test rows ($G_{cs}>10$ W/m$^2$) & 52,731 \\
\bottomrule
\end{tabularx}
\end{table}


Table~\ref{tab:data_quality_summary} is included because the forecasting result depends strongly on what rows are allowed into the evaluation. The active data audit found no missing values in the twelve raw channels, no duplicate timestamps after sorting, and no cadence gaps at the 5-min resolution. During data cleaning, machine-fault periods, hardware-level error states, plant shutdown intervals, and invalid sensor states were removed so that the target represents normal machine operation. Weather-driven disturbances were retained rather than deleting cloudy, humid, rainy, or ramp-heavy intervals. The table distinguishes raw chronological split rows from rows available after lag and rolling features are constructed.

\begin{figure}[!t]
\centering
\includegraphics[width=\columnwidth]{figures/diagrams/data_quality_summary_panel.png}
\caption{Data-quality and descriptive-statistics panel regenerated from the active dataset. The panel summarizes split sizes, missing values, monthly power behavior, and mean cloud-layer levels.}
\label{fig:data_quality_panel}
\end{figure}

\begin{figure}[!t]
\centering
\includegraphics[width=\columnwidth]{figures/results/weather_window_example.png}
\caption{Representative weather-input window used by the notebook diagnostics. The forecast path depends on weather and cloud-layer availability at the target timestamps.}
\label{fig:weather_window}
\end{figure}

Before modeling, timestamps are converted to cyclical time features and chronological blocks. The training interval covers the earliest years, the validation period begins on 2024-01-01, and the final test interval begins on 2024-12-28. This ordering keeps validation and test dates strictly later than the training dates, so model selection and final reporting do not use future rows. The split is shown in Fig.~\ref{fig:split_timeline}.

\begin{figure}[!t]
\centering
\includegraphics[width=\columnwidth]{figures/diagrams/diagram_split_timeline.png}
\caption{Chronological split used for model selection and final testing. The headline results are not based on a random shuffle.}
\label{fig:split_timeline}
\end{figure}
